We test whether Belgian downstream buyers reallocated their input mix in response to the EU ETS at two margins: within a narrowly defined supplier sector (NACE 4-digit), across suppliers; and across NACE 4-digit input categories. On the first margin, the share of the most carbon-exposed ETS supplier in a buyer-cell does not change differentially after 2015 in the larger sample; the most-favorable subset (cells where the supplier's NACE 4-digit accounts for the largest share of the buyer's total input bill) yields a coefficient of $+0.03$ percentage points (s.e.\ $1.39$). On the second margin, the binary regulated $\times$ \textit{Post} specification --- the most direct CMdG-style domestic analog --- yields a coefficient of $-0.003$ (s.e.\ $0.009$), and the year-by-year detrended event study shows uniformly positive coefficients in every clean post-2015 year, pointing in the opposite direction from substitution. A standard Grossman--Krueger emissions decomposition on the same 2005--2022 window shows that within-sector reallocation accounts for essentially zero of the 46\% aggregate emissions decline, consistent with --- though not equivalent to --- the buyer-side null at the within-NACE-4d margin.

We have framed the result conservatively. Three caveats stand: a parallel-trends violation in the within-NACE-4d test that the post-hoc detrended specification handles but does not fully resolve; a small but statistically significant extensive-margin response in which the most-exposed supplier becomes mildly less likely to remain active in the cell; and limited statistical power against modest substitution dispersed across the population at exposure levels below the right tail. With these caveats acknowledged, the joint pattern across margins, sample subsets, and specifications is difficult to reconcile with first-order reallocation in response to Phase~III/IV EU ETS prices in Belgium.

We do not interpret this as a structural statement about the elasticity of substitution between regulated and unregulated suppliers, or as a prediction that buyers \textit{would not} reallocate at higher carbon prices. Whether reallocation would appear at $\sim$2.5$\times$ the realized 2022 EUA price --- where right-tail pair-shock magnitudes would scale to match \citet{peter2023}'s p95 imported-input shock --- is a question we cannot answer with current data, and we leave it open.
